% Generated draft: requires mathematical and source review.
\documentclass[11pt]{article}
\usepackage[margin=25mm]{geometry}
\usepackage{fontspec}
\setmainfont{Latin Modern Roman}
\usepackage{amsmath,amssymb,mathtools,unicode-math}
\setmathfont{Latin Modern Math}
\usepackage{enumitem,needspace,xurl,hyperref}
\setlength{\parindent}{0pt}
\setlength{\parskip}{5pt}
\setlength{\emergencystretch}{3em}
\newcommand{\N}{\mathbb{N}}
\newcommand{\Z}{\mathbb{Z}}
\newcommand{\Q}{\mathbb{Q}}
\newcommand{\R}{\mathbb{R}}
\newcommand{\C}{\mathbb{C}}
\newcommand{\F}{\mathbb{F}}
\newcommand{\A}{\mathbb{A}}
\newcommand{\PP}{\mathbb{P}}
\newcommand{\E}{\mathbb{E}}
\newcommand{\eps}{\varepsilon}
\newcommand{\dd}{\,\mathrm{d}}
\begin{document}
\section*{UnsolvedMath ID: 8800038}
\textbf{Source problem code:} catalogue entry
\subsection{Definitions and mathematical statement}
Let $q$ range over prime powers and let $E/\F_q$ be an elliptic curve, meaning a smooth projective genus-one curve with a specified $\F_q$-rational origin. Prime order means that the finite group $E(\F_q)$ has prime cardinality:
\[
 \#E(\F_q)\text{ is a prime integer}.
\]
It does not mean that $q$ is prime or merely that $E(\F_q)$ contains a point of prime order.

Under an isomorphism-class counting convention, put
\[
 A(q)=\#\{[E]_{\F_q}:\#E(\F_q)\text{ is prime}\}.
\]
The source asks to prove that there are ``sufficiently many'' such curves over every finite field. It does not specify the desired lower bound, asymptotic size, proportion, or quantitative threshold. It also does not choose between counting $\F_q$-isomorphism classes, weighted classes, or curve models; these are not equivalent counting conventions.

The universal field scope is all prime powers $q$, with no source-stated exception for small fields or restriction to prime fields. A precise abundance theorem must supply its count and lower bound and explain whether finitely many exceptions are intended. Merely constructing a special curve for selected orders or fields is not automatically the stated all-fields abundance result. No guessed heuristic count or asymptotic formula is substituted.
\subsection{Short English statement}
Prove an adequately large supply of elliptic curves with prime group order over each finite field. The intended meaning of adequately large remains unspecified.

The supplied summary calls the 2006 framing heuristic, reports restricted complex-multiplication constructions, and connects it with Koblitz/Lang--Trotter questions, but gives no exact abundance theorem. Such fixed-curve or special-order contexts are not silently identified with varying curves over every field. The registry partial-progress label is retained with current primary-literature audit required.
\subsection{Sources}
\begin{itemize}
\item Local UnsolvedMath registry, permanent ID 8800038, source code catalogue entry, ``Prime-order curves over every finite field,'' proposed by Igor Shparlinski. Dataset curation and metadata: Ulam AI and UnsolvedMath Contributors, CC BY 4.0.
\item IPAM, \emph{Workshop on Number Theory and Cryptography --- Open Problems} (2006), Attachment 6224, PDF page 20, ``Holy Grail,'' schedule locator \url{https://www.ipam.ucla.edu/programs/workshops/workshop-i-number-theory-and-cryptography-open-problems/?tab=schedule}; pdf-and-legacy-ppt visual-audit extraction, accessed 2026-07-29.
\item The summary mentions Koblitz, Lang--Trotter, and special complex-multiplication constructions without complete primary references or a quantitative theorem matching this task.
\item Source provenance reports inspection of 17 PDFs and two legacy PPT decks among 19 surviving attachments, excluding repetitions and solved definitions; this is supplied provenance, not a new inspection claimed here.
\item Current primary-literature status audit and redistribution terms NEEDS\_REVIEW for the publicly accessible author/publisher copy, separately from the dataset licence.
\end{itemize}
\end{document}
